\section{Jeux de tests et justification}

\subsection{Interface}
\subsubsection{Problèmes rencontrés}
La réalisation des tests pour l’interface a été plus compliquée que prévu. L’utilisation de TestFX nous a causé des différences de comportement entre les différentes plateformes sur lesquelles nous développons (macOS, Linux, Windows, WSL). Parfois, les tests passaient chez certains, mais échouaient chez d’autres.

Par ailleurs, en ce qui concerne les tests d’acceptation, le comportement de TestFX était encore plus aléatoire d’un test à l’autre, même en utilisant les mêmes fonctions. Suite à ces observations, nous avons décidé d’écrire les scénarios de tests d’acceptation, mais de les exécuter manuellement.

\subsubsection{Choix des tests à effectuer}
Pour la partie interface, nous avons réalisé plusieurs tests. Dans un premier temps, nous avons testé les différents composants graphiques que nous avons réalisés (par exemple la console ou encore la zone de code). Ces tests ont pour but de vérifier que les modifications apportées à ces composants n’entraînent pas de régression dans l’application produite.

Nous avons également réalisé des tests d’acceptation pour chaque user story afin de valider la bonne réalisation de ces user stories et de vérifier que nous n’avons rien oublié.